\chapter*{Úvod}\addcontentsline{toc}{chapter}{Úvod}\markboth{Úvod}{Úvod}
Skauting predstavuje medzinárodné mládežnícke hnutie zamerané na rozvoj osobnosti prostredníctvom praktických aktivít v prírode, spolupráce v kolektíve a neformálneho vzdelávania. Na týchto princípoch funguje dobrovoľnícka organizácia Slovenský skauting, ktorá na Slovensku zabezpečuje programové aktivity pre deti a mládež.

Program organizácie je realizovaný najmä prostredníctvom pravidelných stretnutí družín, ktoré predstavujú základnú organizačnú jednotku skautingu. Družina je spravidla tvorená skupinou približne dvanástich členov vedených radcom. Radca je zodpovedný za prípravu programu stretnutí a koordináciu aktivít družiny.

Pri plánovaní programu však radcovia často čelia náročnej úlohe zosúladiť spoločné aktivity družiny s individuálnym napredovaním jednotlivých členov v rámci skautského programu.

Napriek existencii programovej štruktúry organizácie Slovenský skauting môžu radcovia a vedúci pri plánovaní aktivít a sledovaní pokroku členov narážať na viacero praktických problémov. Programová ponuka je síce dostupná prostredníctvom oficiálnych webových stránok organizácie a tlačených príručiek, ktoré obsahujú informácie o stupňoch napredovania, odborkách a ďalších programových prvkoch, ich používanie pri každodennej práci s družinou však je často zdĺhavé a nepraktické.

Zároveň existuje veľké množstvo digitálnych nástrojov využívaných na komunikáciu, organizáciu aktivít alebo evidenciu úloh. Používanie viacerých oddelených nástrojov však môže viesť k roztriešteniu informácií a zvyšovať organizačnú záťaž vedúcich a radcov. Táto situácia vytvára priestor pre komplexné riešenie, ktoré by integrovalo prehľadné sprístupnenie programovej ponuky s jednoduchým nástrojom na sledovanie pokroku členov a podporu organizácie aktivít družiny.

\section*{Cieľ práce}\addcontentsline{toc2}{section}{Cieľ práce}\markboth{Cieľ práce}{Cieľ práce}
Cieľom tejto diplomovej práce je návrh, implementácia a otestovanie aplikácie určenej pre členov organizácie Slovenský skauting. Aplikácia má prehľadnejšie sprístupniť programovú ponuku organizácie jej členom a znížiť organizačnú záťaž pri sledovaní individuálneho pokroku členov zo strany radcov a vedúcich.

Aby aplikácia naplnila svoj účel, musí byť intuitívna pre široké vekové spektrum používateľov. Súčasťou práce je preto dôkladný návrh používateľského rozhrania (UI/UX), samotná implementácia systému, ktorý zjednotí zobrazovanie programových prvkov so správou družinovej agendy, a následné používateľské testovanie vytvoreného riešenia.

Práca je rozdelená do šiestich hlavných kapitol. Prvá kapitola sa venuje teoretickému základu, kde definujeme princípy organizácie Slovenský skauting, analyzujeme existujúce digitálne riešenia a popisujeme technologické východiská z oblasti softvérového inžinierstva a tvorby používateľských rozhraní. Druhá kapitola obsahuje analýzu cieľovej skupiny a definovanie funkčných aj nefunkčných požiadaviek na aplikáciu. V tretej kapitole popisujeme proces návrhu, od doménového modelu a architektúry až po tvorbu a testovanie prototypov nízkej a vysokej vernosti. Štvrtá kapitola je zameraná na samotnú implementáciu v prostredí Flutter a proces získavania dát. Piata kapitola vyhodnocuje úspešnosť riešenia prostredníctvom používateľského testovania, heuristickej analýzy a automatizovaných testov. V závere sumarizujeme dosiahnuté výsledky a navrhujeme možnosti pre budúci vývoj aplikácie.